\section{Physical clearance certificates and complete-source localization}
\label{sec:gv-source}
The endpoint defect in \eqref{eq:gv-endpoint-bound} is available for
every original regular trajectory, including those whose geometric
word minimum is not a true critical orbit. We quantify this distinction
rather than discarding that source. The final proposition then inserts
the proved enclosure into the existing exact source formula, at the
original fixed radius and roof.

\begin{theorem}[A quantitative certificate relative to every minimizing contact]
\label{thm:gv-clearance-certificate}
Let $Q$ be an arbitrary true regular record and let $E=E_R(Q)$.
For every minimizer $Y$ defining $\ell_w(R)$ set
\[
 \delta_0=\delta_m=2\sqrt{2RE},\qquad
 \delta_i=\sqrt{RE/c_i}\quad(0<i<m).
\]
Then $|Y_i-Q_i|\le\delta_i$ for every $i$. Consequently, for every
unvisited center $z$ on every original flight, without a lattice
truncation,
\begin{equation}\label{eq:gv-clearance-certificate}
 \operatorname{dist}(z,[Y_i,Y_{i+1}])-R
 \ge\operatorname{dist}(z,[Q_i,Q_{i+1}])-R
                    -\max(\delta_i,\delta_{i+1}).
\end{equation}
There is a sufficient \emph{proved} physicality certificate: if
\begin{align}
 \Delta_i&>\max(\delta_i,\delta_{i+1})\quad(0\le i<m),
                                                     \label{eq:gv-clearance-test}\\
 c_i&>(\delta_{i-1}+2\delta_i+\delta_{i+1})/\tau_-
                      \quad(0<i<m),                 \label{eq:gv-incidence-test}
\end{align}
then every minimizer $Y$ is the same regular physical normal-endpoint
record of the original word. No section label is thereby transferred
from $Q$ to $Y$. If these tests fail, $Q$ stays in the original source;
its minimum is not declared physical.
\end{theorem}
\begin{proof}
Insert the actual velocity vector of $Q$ and a minimizing $Y$ into
\eqref{eq:gv-gap}. Weak duality and \eqref{eq:gv-endpoint-exact} give
$0\le S(Y)-D_{w,R}(V)\le E$. Thus each nonnegative contact term is
at most $E$. At an interior contact it equals
$2c_i\{R-n_i\cdot(Y_i-O_i)\}\ge(c_i/R)|Y_i-Q_i|^2$,
which proves the interior bound.

At departure the supporting disk point for $-v_0$ is $O_0+Rv_0$;
the endpoint contact slack bounds its squared distance from $Y_0$
by $2RE$. Also
$|Q_0-(O_0+Rv_0)|^2=2R^2(1-c_0)\le2RE$.
The triangle inequality gives $|Y_0-Q_0|\le2\sqrt{2RE}$.
At arrival use $O_m-Rv_{m-1}$ and the same argument.
For a common affine interpolation parameter the distance between two
corresponding segment points is at most
$\max(\delta_i,\delta_{i+1})$. Both directed Hausdorff distances
satisfy this bound. Distance to any fixed $z$ is 1-Lipschitz for
Hausdorff distance; this proves \eqref{eq:gv-clearance-certificate}
for all lattice centers simultaneously.

It remains to justify physicality, not merely small displacement.
The edges of $Y$ have length at least $\tau_-=3/50$, because
successive original disks are disjoint. If $\widehat v_i$ is their
unit direction, the normalization inequality
$|u/|u|-v/|v||\le2|u-v|/\min(|u|,|v|)$ gives
\[
 |\widehat v_i-v_i|\le2(\delta_i+\delta_{i+1})/\tau_-.
\]
Hence the $i$th interior gradient of $S$ at $Y$ has norm at least
$2c_i-2(\delta_{i-1}+2\delta_i+\delta_{i+1})/\tau_->0$.
By the minimizing condition in Theorem~\ref{thm:gv-duality}, $Y_i$
lies on its circle and that gradient is a negative multiple of the
outward normal. Equality of the two adjacent unit speeds then makes
the incoming normal component strictly negative and the outgoing one
strictly positive, with equal tangential components: this is exactly
the regular exterior specular law. At the endpoints the gradients
have norm one and give normal departure and arrival. The strict
third-disk inequalities \eqref{eq:gv-clearance-test} exclude every
unlisted collision; the signs just proved exclude another intersection
with either incident disk inside its flight. Thus $Y$ is a true regular
first-hit record. Corollary~\ref{cor:gv-critical} gives uniqueness.
Nothing in the proof keeps section visits constant along a competitor
or along a homotopy; the original labels of $Q$ have not changed.
\end{proof}

\begin{proposition}[A deterministic positive junction cost]
\label{prop:gv-junction}
Partition a word into consecutive nonempty blocks at any chosen
indices. For each block choose a minimizing unit-edge dual vector as
in Theorem~\ref{thm:gv-duality}. Denote the last vector before a
junction by $v^-$ and the first after it by $v^+$. Then
\begin{equation}\label{eq:gv-junction}
 \ell_w(R)\ge\sum_{\text{blocks }b}\ell_b(R)
                  +R\sum_{\text{junctions}}(2-|v^--v^+|).
\end{equation}
Every junction term is nonnegative. If both block minima at a junction
are physical critical records with terminal/departure normals
$n^-,n^+$, the corresponding term is $R(2-|n^-+n^+|)$.
The comparison does not assert physical realizability of any block
minimum and does not factor any trajectory probability.
\end{proposition}
\begin{proof}
Concatenate the chosen dual flight vectors. They form a feasible dual
vector for the whole word. The center term $C_w$ is the sum of the
block center terms. At a junction the two separate endpoint penalties
are $R|v^-|+R|v^+|=2R$, whereas the penalty of the concatenated dual
vector is $R|v^--v^+|$. All other penalties agree. Substitution into
\eqref{eq:gv-duality} proves \eqref{eq:gv-junction}, including each
junction once. In the physical critical case $v^-=-n^-$ and
$v^+=n^+$, giving the asserted normal expression. This argument is
an exact vector comparison, not a renewal or Markov decomposition.
\end{proof}

\begin{proposition}[Word localization of the unchanged complete roof source]
\label{prop:gv-source}
Fix the original $m,R,n,k$ and any nonnegative Borel trajectory weight
$h$. Let $x_w(t,\vartheta)$ be exactly the partial physical inverse
in \eqref{eq:v111-roof-sum}, with all original event cuts. That formula
admits the additional exact mask
\begin{equation}\label{eq:gv-source-mask}
 \mathbf1_{\{\ell_w(R)\le t\le\ell_w(R)+E_R(x_w(t,\vartheta))\}}
\end{equation}
in every existing summand, without changing either side. In particular
only words with $\ell_w(R)\in[t-2R,t]$ can contribute. This assertion
requires no physical critical center, clearance guard, pressure
hypothesis or incidence lower bound for the regular source.

For clarity about what a pressure consequence does and does not add,
define the deterministic localized physical-word pressure
\[
 P_{m,R}(J)=\sum_{\substack{w\in\mathcal W_m(R)\\\ell_w(R)\in J}}
                       \sup_{[w]_R\cap M_R'}j_m.
\]
Under \emph{exactly} Input~\ref{input:v112-geometric-pressure} and
the inherited roof--angle formula, for $0\le h\le1$,
\begin{equation}\label{eq:gv-localized-pressure}
 \frac{dL_\#(h\mathbf1_{\Pi_{n,k,m,R}}\nu_R^*)}{dt}(t)
 \le\frac{C_0}{2c_*\lambda_\circ}\,P_{m,R}([t-2R,t])
 \quad\text{for almost every }t.
\end{equation}
This is an unscaled comparison, not a new estimate on that localized
pressure and not a proof of the original normalized shortage.
\end{proposition}
\begin{proof}
Theorem~\ref{thm:gv-endpoint} applies to every existing physical
inverse point, so its mask is one. Apply the nonnegative source
exhaustion of Proposition~\ref{prop:v112-exhaustion}; the mask is
Borel, and no test or boundary is differentiated. This proves the
exact identity with the mask, for extended weights as well.
For the last bound use $\chi\ge\lambda_\circ$ and
$\omega\le C_0c_mj_m\le C_0j_m$ at the actual point, and dominate
$h$ and the actual event only on the upper side by one. Global
physical roof--angle uniqueness gives at most one point per word.
Each word retained by the mask lies in the indicated geometric
interval, so its stable weight is bounded by the same-word supremum
at the \emph{same} $R$. Integration over the angle of length $2\pi$
gives the factor in \eqref{eq:gv-localized-pressure}. The section
factor $c_*^{-1}$ occurs once, exactly as in the inherited formula.
The supremum of this unlocalized pressure is the retained $Q_*$;
no decay in $m$ for the localized expression has been inferred.
\end{proof}

\subsection{The original positive shortage is still the closing condition}
The new results remove the need for a physical critical anchor in the
geometric word enclosure; certify an actual-versus-unphysical minimizing
configuration using nonnegative gaps; give a positive junction defect;
and compare critical roofs across disconnected physical radius domains.
They do not identify the auxiliary $\ell_w$ with a physical clock when
the certificate fails. In particular \eqref{eq:gv-source-mask} does not
replace the observed microscopic window by $[t-2R,t]$. The observed
quantity is still the density at the original $t$ with the original
exact labels; the interval is solely an upper-support enclosure in a
separate deterministic word statistic.

No new theorem here supplies a normalized roof-height estimate for the
remaining positive source or enough nongrazing comparison capacity.
The raw target is still \eqref{eq:v76-target}, with the very same minimal
ordered closure condition \eqref{eq:v112-remaining}. It includes the
low-load incidence-only source, first-order unvisited clearance,
noncritical singular charts, one-sided paired/section traces and all
original arithmetic/denominator requirements. The finite source sum is
taken before $\sup_R\esssup_t$ and before the collision limsup; all
microscopic widths remain fixed until that limsup, and $B\to\infty$
is outermost. Neither an auxiliary convex value, a positive curvature
identity, a finite diagnostic, nor the radius-integrated coarea pays
that still unproved limit automatically.
